\chapter{Concluzii}

\section{Rezultate obținute}

Lucrarea de licență a avut ca obiectiv dezvoltarea \textbf{aplicației web pentru organizarea și gestionarea evenimentelor}, orientate întâi spre conferințe, meeting-uri, concerte și evenimente culturale, cu suport complementar pentru competiții sportive. Obiectivele propuse au fost atinse:

\begin{itemize}
	\item Aplicație full-stack funcțională: React + Express + PostgreSQL.
	\item Autentificare completă: înregistrare cu verificare email, login JWT, resetare parolă, profil, ștergere cont.
	\item Listare și filtrare evenimente (nume, județ, dată).
	\item Plăți Stripe Checkout în RON; bilete nominale și ,,Biletele mele``.
	\item Panou admin: CRUD evenimente cu imagini, check-in, gestionare admini.
	\item Documentare: arhitectură, ERD, UML, API REST, scenarii pe tipuri de evenimente.
\end{itemize}

Soluția demonstrează integrarea unor tehnologii moderne într-un produs coerent, adaptat contextului românesc (județe, RON, email Brevo).


\section{Limitări}

\begin{itemize}
	\item Lipsă teste automate (unit/integration) și pipeline CI/CD.
	\item Promovarea primului administrator necesită script SQL sau \texttt{promote-admin}.
	\item Emailurile depind de configurarea SMTP; în development, codurile pot apărea în consola backend.
	\item Confirmarea plății se bazează pe apel frontend \texttt{confirm-payment}; webhook Stripe ar îmbunătăți robustețea.
	\item Un singur tip de preț per eveniment — fără categorii VIP/standard sau locuri numerotate.
\end{itemize}

\section{Dezvoltări viitoare}

Dezvoltarea aplicației poate fi continuată prin abordarea limitărilor identificate și adăugarea de noi funcționalități:

\begin{itemize}
	\item Categorii de bilete (conferință: early bird / regular; concert: VIP).
	\item Export participanți CSV din panoul admin.
	\item Notificări email la cumpărare bilet reușită.
	\item Paginare evenimente și caching.
	\item Webhook Stripe pentru confirmare server-side.
	\item Teste automate și deploy containerizat (Docker).
	\item Aplicație mobilă sau PWA pentru check-in offline.
	\item Module opționale pentru evenimente sportive (echipe, program).
\end{itemize}

Aplicația reprezintă o bază solidă pentru extindere într-un produs destinat organizatorilor de conferințe, concerte și meeting-uri, cu accent pe simplitate, securitate și experiența utilizatorului.